\documentclass[12pt,spanish,letterpaper,oneside]{article}
\def\documenttitle {Antecedentes sociales y contexto pol\'itico de los derechos humanos en M\'exico}
\def\documentsubtitle {Continuidades, transformaciones y l\'imites}
\def\documentsubject {Cuadro comparativo}
\def\documentauthor {Martin Jonathan de la Cruz}
\def\coursename {Antecedentes de los derechos humanos}
\def\coursecode {005}
\def\universityname {Universidad Abierta y a Distancia de M\'exico}
\def\universityfaculty {Licenciatura en Derecho}
\def\universitydepartment {Antecedentes de los derechos humanos}
\def\universitydepartmentimage {departamentos/UnADM}
\def\universitydepartmentimagecfg {height=1.57cm}
\def\universitylocation {Ciudad de M\'exico}
\def\authortable {\begin{tabular}{ll}
	\textbf{Alumno:} & Martin Jonathan de la Cruz \\
	Matricula: & ES2611202040 \\
	Figura docente: & Daniel Alexis Lozano Ortega \\
	Semestre/Bloque: & 2 / 2 \\
	Actividad oficial: & 5 \\
	Semana: & 4 \\
	Fecha: & 22 de octubre de 2026
\end{tabular}}
\usepackage{iftex}
\ifLuaTeX
	\let\pdfpagewidth\pagewidth
	\let\pdfpageheight\pageheight
\fi
\input{template}
\usepackage{fontspec}
\setmainfont{arial.ttf}[Path=C:/Windows/Fonts/,BoldFont=arialbd.ttf,ItalicFont=ariali.ttf,BoldItalicFont=arialbi.ttf]
\setsansfont{arial.ttf}[Path=C:/Windows/Fonts/,BoldFont=arialbd.ttf,ItalicFont=ariali.ttf,BoldItalicFont=arialbi.ttf]
\renewcommand{\familydefault}{\sfdefault}
\usepackage{pdflscape}
\usepackage{array,tabularx,colortbl}
\usepackage{tikz}
\usepackage{eso-pic}
\setcitestyle{authoryear,open={(},close={)}}
\newcommand{\insertcoverwatermark}{%
	\AddToShipoutPictureBG*{%
		\begin{tikzpicture}[remember picture,overlay]
			\node[opacity=0.16,inner sep=0pt] at (current page.center)
				{\includegraphics[width=0.72\paperwidth]{img/departamentos/UnADM.pdf}};
		\end{tikzpicture}%
	}%
}
\definecolor{cabecera}{RGB}{225,237,232}
\renewcommand{\tabularxcolumn}[1]{p{#1}}
\newcolumntype{Y}{>{\raggedright\arraybackslash}X}
\setlength{\tabcolsep}{5pt}
\setlength{\emergencystretch}{3em}
\onehalfspacing
\newcommand{\encabezado}[2]{\rowcolor{cabecera}\textbf{Categor\'ia}&\textbf{#1}&\textbf{#2}\\\hline}
\usepackage{longtable}
\newlength{\categorywidth}
\newlength{\periodwidth}
\newlength{\lastperiodwidth}
\begin{document}
\insertcoverwatermark
\templatePortrait
\templatePagecfg
\onehalfspacing
\begin{abstractd}
Los derechos humanos en M\'exico se comprenden mediante el contraste entre
formas de organizaci\'on social, legitimidad del poder y mecanismos de
protecci\'on. El recorrido desde las sociedades prehisp\'anicas hasta el
paradigma constitucional contempor\'aneo distingue deberes comunitarios,
tutelas coloniales, libertades individuales y derechos sociales. La
comparaci\'on de ocho etapas mediante categor\'ias comunes muestra que la
existencia de normas no equivale a igualdad universal ni a cumplimiento
efectivo. La continuidad de conflictos y exclusiones impide interpretar
la historia como un progreso autom\'atico; el criterio decisivo es c\'omo
se limitan los abusos y se hacen exigibles las obligaciones de protecci\'on.
\end{abstractd}
\templateIndex
\templateFinalcfg
\newgeometry{margin=2.5cm}
\fontsize{12}{14.5}\selectfont
\onehalfspacing

\section{Introducci\'on}
El desarrollo de los derechos humanos en M\'exico comprende conflictos
sociales, transformaciones del poder y cambios en las formas de protecci\'on
jur\'idica. No constituye una secuencia de avances inevitables: el
reconocimiento de derechos puede coexistir con exclusiones y dificultades
para ejercerlos. La historia de su concepci\'on moderna muestra avances y
retrocesos, mientras que la organizaci\'on constitucional permite distinguir
la proclamaci\'on de libertades de las instituciones que limitan el poder
\citep[p.~11]{alvarez2009}\citep[pp.~1--2]{arteaga}.

El recorrido hist\'orico comprende ocho etapas: antecedentes prehisp\'anicos,
Conquista, Virreinato, Ilustraci\'on, Independencia, constitucionalismo liberal,
Revoluci\'on y evoluci\'on contempor\'anea. En las primeras etapas se examinan
formas de organizaci\'on y controversias sobre la protecci\'on de las personas,
sin atribuirles el concepto contempor\'aneo de derechos universales. Las
declaraciones liberales, las reivindicaciones sociales y las garant\'ias actuales
se comparan atendiendo tanto a sus aportaciones como a sus l\'imites.

\section{Dignidad, orden social y transformaciones de la protecci\'on jur\'idica}
\subsection{Fundamento, reconocimiento y garant\'ias}
La existencia de normas no basta para afirmar que una sociedad reconoce
derechos humanos universales. Es necesario distinguir tres planos: las
razones que justifican proteger a las personas, las disposiciones que
reconocen facultades y los mecanismos que permiten exigir su respeto.
Una constituci\'on organiza el poder y garantiza libertades, pero su estructura
normativa no demuestra por s\'i sola el cumplimiento de sus disposiciones
\citep[pp.~1--2]{arteaga}. La historia de los derechos muestra, adem\'as,
avances y retrocesos; por ello, una secuencia cronol\'ogica no debe leerse
como progreso inevitable \citep[p.~11]{alvarez2009}.

En el estudio de los mexicas, el calpulli, el gobierno y el tributo expresan
formas de organizaci\'on comunitaria y relaciones jer\'arquicas. La conservaci\'on
de autoridades e idiomas en pueblos dominados no equivale a igualdad de
posiciones. Estos antecedentes permiten estudiar normas y deberes sin
proyectar sobre toda Mesoam\'erica una sola organizaci\'on ni atribuirle el
concepto constitucional contempor\'aneo de derechos humanos
\citep[cap.~1]{brom2017}.

La encomienda muestra una tensi\'on distinta: el deber declarado de proteger
y evangelizar coexist\'ia con la exigencia de tributo y servicios. Brom
describe esta organizaci\'on como una relaci\'on de dominio y explotaci\'on;
no como propiedad de las tierras ni como prueba de tutela efectiva.
Las rep\'ublicas de indios y la tierra comunal, junto a las autoridades
virreinales, muestran que pod\'ian existir protecciones particulares en un
orden social desigual \citep[caps.~5--6]{brom2017}. La presencia de una
instituci\'on protectora debe examinarse, por tanto, junto con la posici\'on
de quienes acud\'ian a ella y los l\'imites de su actuaci\'on.

Las declaraciones liberales desplazan el problema hacia la libertad, la
igualdad jur\'idica y la legitimidad del gobierno representativo
\citep[pp.~11--13]{alvarez2009}. En la insurgencia mexicana, la soberan\'ia
popular, la divisi\'on de poderes y la prohibici\'on de esclavitud y castas
se articulan con la exigencia de moderar la opulencia y la indigencia
\citep[nums.~5, 12 y 15]{morelos1813}. Esa combinaci\'on impide reducir
la Independencia a un cambio de autoridad: tambi\'en se disputaba el
contenido social del nuevo orden. No significa que todos los proyectos
hayan tenido vigencia inmediata.

El constitucionalismo social agrega condiciones materiales de educaci\'on,
tierra y trabajo a la protecci\'on de libertades. Los contenidos sociales
y la organizaci\'on del poder permanecen relacionados en la Constituci\'on
\citep[p.~2]{arteaga}. El texto vigente de los art\'iculos 3, 27 y 123
incluye reformas posteriores y no debe reproducirse como si fuera
\'integramente la redacci\'on de 1917 \citep{cpeum2026}.

En el paradigma actual, el art\'iculo 1 vincula derechos constitucionales
y convencionales con obligaciones de todas las autoridades. Universalidad,
interdependencia, indivisibilidad y progresividad orientan su protecci\'on;
el principio pro persona exige favorecer la protecci\'on m\'as amplia
\citep[art.~1]{cpeum2026}. Los organismos no jurisdiccionales constituyen
otra v\'ia de defensa, pero sus recomendaciones no equivalen a sentencias
\citep[pp.~47--48]{alvarez2009}. Estas diferencias permiten comparar las
etapas por su organizaci\'on social, legitimidad del poder, contenido de
los derechos y eficacia de los instrumentos de protecci\'on.

% Ficha oficial web 76: conceptos en columnas y categorias comunes en filas.
% Cuatro comparaciones cualitativas conservan las ocho etapas de Unidad 3.
% Tabla continua compacta: dos paginas horizontales, cuatro periodos por pagina.
\begin{landscape}
\pagestyle{empty}
\singlespacing
\setlength{\parindent}{0pt}
\setlength{\parskip}{0pt}
\setlength{\tabcolsep}{4pt}
\setlength{\extrarowheight}{0pt}
\renewcommand{\arraystretch}{1}
\fontsize{9}{10.5}\selectfont
\captionsetup{font=footnotesize}
\setlength{\LTleft}{0pt}
\setlength{\LTright}{0pt}
\setlength{\categorywidth}{\dimexpr0.16\linewidth-2\tabcolsep\relax}
\setlength{\periodwidth}{\dimexpr0.42\linewidth-2\tabcolsep\relax}
\setlength{\lastperiodwidth}{\dimexpr0.42\linewidth-2\tabcolsep-4\arrayrulewidth\relax}
\setlength{\LTcapwidth}{\linewidth}
\captionof{table}{Cuadro comparativo de los antecedentes sociales y politicos de los derechos humanos en Mexico}\label{tab:antecedentes-comparativo}
\begin{longtable}{|>{\raggedright\arraybackslash}p{\categorywidth}|>{\raggedright\arraybackslash}p{\periodwidth}|>{\raggedright\arraybackslash}p{\lastperiodwidth}|}
\hline\encabezado{Antecedentes prehisp\'anicos}{Conquista y defensa ind\'igena}
\endfirsthead
\multicolumn{3}{c}{\footnotesize Cuadro comparativo: antecedentes sociales y politicos (continuacion)}\\
\hline
\endhead
\multicolumn{3}{r}{\footnotesize Contin\'ua en la siguiente p\'agina}\\
\endfoot
\endlastfoot
Contexto hist\'orico & Antes de 1521; sociedades diversas. La muestra mexica no representa toda Mesoam\'erica. & Siglo XVI; conquista y establecimiento del dominio espa\~nol. \\\hline
Antecedentes sociales & Calpulli, actividades productivas y educaci\'on diferenciada por clase. & Exigencia de tributo y servicios mediante la encomienda; alteraci\'on de formas de vida. \\\hline
Contexto pol\'itico & Hueytlatoani y tlatocan; pueblos tributarios conservaban autoridades e idiomas. & Monarqu\'ia y evangelizaci\'on; concesi\'on y limitaci\'on de encomiendas. \\\hline
Acontecimientos & Organizaci\'on comunitaria y tributaria; conflictos entre dominadores y pueblos sometidos. & Ca\'ida de Tenochtitlan; expansi\'on territorial; defensa de ind\'igenas por Las Casas y Quiroga. \\\hline
Derechos o principios & Deberes comunitarios y orden social; no equivalen a derechos universales actuales. & Protecci\'on y evangelizaci\'on declaradas; cr\'iticas al sometimiento, sin igualdad efectiva general. \\\hline
Aportaciones y l\'imites & Permite estudiar diversidad y jerarqu\'ias sin idealizar la comunidad \citep[cap.~1]{brom2017}. & Expone la distancia entre tutela declarada y explotaci\'on. Encomienda no equivale a propiedad de tierras \citep[caps.~4--5]{brom2017}. \\\hline
Semejanza y diferencia & Ambos periodos presentan tributo y subordinaci\'on; la organizaci\'on mexica conservaba formas comunitarias propias. & Cambian los beneficiarios y la legitimaci\'on del dominio: monarqu\'ia y evangelizaci\'on. No es una simple continuidad institucional. \\\hline
\encabezado{Virreinato de Nueva Espa\~na}{Ilustraci\'on y revoluciones liberales}
Contexto hist\'orico & Siglos XVI--principios del XIX; organizaci\'on colonial. & Siglos XVII--XVIII; referentes europeos y norteamericanos. \\\hline
Antecedentes sociales & Peninsulares, criollos, ind\'igenas y castas; tierra comunal bajo control colonial. & Cuestionamiento de privilegios e impuestos; demandas de representaci\'on. \\\hline
Contexto pol\'itico & Rey y Consejo de Indias; virrey y Audiencia con funciones judiciales y de control. & Cr\'itica del absolutismo; soberan\'ia, representaci\'on y divisi\'on de poderes. \\\hline
Acontecimientos & R\'egimen virreinal desde 1535; rep\'ublicas de indios y espa\~noles. & Independencia estadounidense (1776) y Declaraci\'on francesa (1789). \\\hline
Derechos o principios & Tutelas y jurisdicciones diferenciadas por condici\'on; no ciudadan\'ia igualitaria. & Libertad, igualdad jur\'idica, propiedad, seguridad y resistencia a la opresi\'on. \\\hline
Aportaciones y l\'imites & Instituciones protectoras coexist\'ian con jerarqu\'ias y subordinaci\'on \citep[caps.~5--6]{brom2017}. & Lenguaje de derechos y legitimidad constitucional; proclamaci\'on no elimina exclusiones \citep[pp.~11--13]{alvarez2009}. \\\hline
Semejanza y diferencia & Ambos ordenan relaciones entre autoridad y personas; el orden colonial diferencia posiciones jur\'idicas. & El ideal liberal cuestiona privilegios y propone igualdad jur\'idica; su alcance efectivo debe contrastarse, no presumirse. \\\hline
\pagebreak
\encabezado{Independencia}{Constitucionalismo liberal del siglo XIX}
Contexto hist\'orico & 1810--1821; ruptura colonial y proyectos nacionales. & Organizaci\'on nacional tras 1821; Constituci\'on de 1857. \\\hline
Antecedentes sociales & Demandas populares y conflicto criollo con privilegios peninsulares \citep[p.~73]{delgado2010}. & Conflictos sobre privilegios y libertades; persistencia de desigualdad material. \\\hline
Contexto pol\'itico & Disputa de soberan\'ia y legitimidad; independencia y representaci\'on. & Proyectos de Estado y supremac\'ia constitucional; l\'imites al gobierno. \\\hline
Acontecimientos & Sentimientos de la Naci\'on (1813), Apatzing\'an (1814) y consumaci\'on de 1821. & Textos constitucionales y reformas; 1857 como antecedente del texto reformado en 1917. \\\hline
Derechos o principios & Soberan\'ia popular, divisi\'on de poderes, proscripci\'on de esclavitud y castas; moderaci\'on de desigualdades. & Libertades individuales, legalidad y protecci\'on frente a actos de autoridad. \\\hline
Aportaciones y l\'imites & Vincula emancipaci\'on con transformaci\'on social; propuesta no equivale a vigencia efectiva \citep[nums.~5, 12 y 15]{morelos1813}. & Organiza el poder y protege libertades; igualdad formal no resuelve desigualdad econ\'omica \citep[pp.~1--2, 7]{arteaga}. \\\hline
Semejanza y diferencia & Ambos cuestionan el poder arbitrario. La insurgencia disputa la ruptura colonial y el nuevo sujeto soberano. & El constitucionalismo desarrolla la organizaci\'on y los l\'imites jur\'idicos del Estado independiente. \\\hline
\encabezado{Revoluci\'on y Constituci\'on de 1917}{Evoluci\'on contempor\'anea}
Contexto hist\'orico & Demandas revolucionarias y constitucionalismo social. & Sistema p\'ublico de defensa y derechos humanos constitucionales. \\\hline
Antecedentes sociales & Demandas agrarias, laborales y educativas frente a desigualdades materiales. & Organizaci\'on ciudadana, democracia, no discriminaci\'on y atenci\'on a violaciones. \\\hline
Contexto pol\'itico & Poder constituyente y ampliaci\'on de la intervenci\'on social del Estado. & Organismos de defensa y articulaci\'on de Constituci\'on y tratados. \\\hline
Acontecimientos & Promulgaci\'on de 1917; articulaci\'on de libertades y contenidos sociales. & CNDH (1990), rango constitucional (1992), autonom\'ia (1999); reforma de 2011. \\\hline
Derechos o principios & Educaci\'on, tierra y trabajo; referentes en los art\'iculos 3, 27 y 123. & Pro persona, interdependencia y progresividad; deberes de promover, respetar, proteger y garantizar. \\\hline
Aportaciones y l\'imites & Incorpora condiciones materiales. El texto vigente incluye reformas posteriores a 1917 \citep[p.~2]{arteaga}; \citep[arts.~3, 27 y 123]{cpeum2026}. & Articula obligaciones y defensa no jurisdiccional. Recomendaciones no equivalen a sentencias \citep[pp.~47--48]{alvarez2009}; \citep[arts.~1 y 102 B]{cpeum2026}. \\\hline
Semejanza y diferencia & Ambos exigen acci\'on estatal; el constitucionalismo social incorpora demandas materiales al texto constitucional. & El paradigma actual articula derechos y tratados y precisa deberes de protecci\'on. No sustituye los derechos sociales. \\\hline
\end{longtable}
\end{landscape}
\clearpage
\pagestyle{plain}
\onehalfspacing

\subsection{Semejanzas y diferencias}
Las ocho etapas relacionan formas de organizaci\'on social con el ejercicio
del poder. Su semejanza principal es que muestran tensiones entre las
reglas proclamadas y las condiciones reales de protecci\'on. No son,
sin embargo, equivalentes: pertenecer a una comunidad o contar con una
protecci\'on particular en el orden colonial no significa ser titular de
derechos universales en el sentido constitucional actual.

La Ilustraci\'on y la Independencia comparten el cuestionamiento de la
legitimidad del poder, pero difieren en su escala: las declaraciones liberales
formulan principios generales, mientras que la insurgencia mexicana los
vincula con emancipaci\'on y construcci\'on nacional. El constitucionalismo
liberal prioriza libertades y l\'imites jur\'idicos; el social incorpora demandas
agrarias y laborales. Ambos contenidos permanecen en la Constituci\'on y
no se sustituyen entre s\'i \citep[pp.~1--2]{arteaga}.

La evoluci\'on contempor\'anea agrega organismos de defensa y una
articulaci\'on expl\'icita con tratados internacionales. Su diferencia no es
solo ampliar el cat\'alogo: el art\'iculo 1 exige obligaciones a todas las
autoridades y ordena interpretar favoreciendo la protecci\'on m\'as amplia
\citep[art.~1]{cpeum2026}. La distancia entre reconocimiento y cumplimiento
sigue siendo el punto com\'un que debe evaluarse en cada periodo.

\clearpage
\section{Conclusiones}
El estudio comparado permite comprender que los derechos actuales no
proceden de una sola declaraci\'on ni de un progreso lineal. Sus antecedentes
muestran conflictos sobre dignidad, legitimidad del gobierno, libertades,
condiciones materiales y mecanismos de defensa. Reconocer esos cambios
impide confundir instituciones antiguas con derechos universales o atribuir
a la promulgaci\'on de una norma una eficacia autom\'atica.

Desde esta perspectiva, estudiar el pasado resulta relevante para valorar
el sistema actual: permite preguntar qui\'en recibe protecci\'on, qu\'e autoridad
est\'a obligada y qu\'e medio permite reclamarla. La importancia de un
acontecimiento no depende solo del principio que enuncia, sino tambi\'en
de las exclusiones que combate y de los instrumentos que deja disponibles.
Ese criterio ofrece una base argumentada para examinar la protecci\'on
efectiva de la dignidad humana en M\'exico.\footnote{GitHub Copilot apoy\'o
la organizaci\'on conceptual y el contraste documental; la postura final
requiere revisi\'on personal del estudiante.}

\clearpage
\begin{thebibliography}{6}
\bibitem[\'Alvarez Icaza Longoria(2009)]{alvarez2009}
\'Alvarez Icaza Longoria, E. (2009). \emph{Los derechos humanos en M\'exico}. Nostra Ediciones.
\bibitem[Arteaga Nava(2014)]{arteaga}
Arteaga Nava, E. (2014). \emph{Derecho constitucional} (4.\textsuperscript{a} ed., edici\'on digital). Oxford University Press.
\bibitem[Brom(2017)]{brom2017}
Brom, J. (2017). \emph{Esbozo de historia de M\'exico} (1.\textsuperscript{a} ed. digital). Penguin Random House Grupo Editorial.
\bibitem[C\'amara de Diputados(2026)]{cpeum2026}
C\'amara de Diputados del H. Congreso de la Uni\'on. (2026).
\emph{Constituci\'on Pol\'itica de los Estados Unidos Mexicanos}
(\'ultima reforma DOF 7 de octubre de 2026).
\url{https://www.diputados.gob.mx/LeyesBiblio/pdf/CPEUM.pdf}
\bibitem[Delgado de Cant\'u(2010)]{delgado2010}
Delgado de Cant\'u, G. M. (2010). \emph{Historia universal: De la era de las revoluciones al mundo globalizado} (3.\textsuperscript{a} ed.). Pearson Educaci\'on.
\bibitem[Morelos y Pav\'on(1813/2021)]{morelos1813}
Morelos y Pav\'on, J. M. (2021). Sentimientos de la Naci\'on.
En M. \'A. Fern\'andez Delgado, \emph{Los Sentimientos de la Naci\'on de Jos\'e Mar\'ia Morelos}.
Instituto Nacional de Estudios Hist\'oricos de las Revoluciones de M\'exico.
(Obra original publicada en 1813).
\url{https://inehrm.gob.mx/recursos/Libros/Los_Sentimientos_de_la_Nacion_de_Jose_Maria_Morelos.pdf}
\end{thebibliography}
\end{document}